\documentclass[../../proceedings/main.tex]{subfiles}
\begin{document}
\begin{refsection}[\subfix{references.bib}]
\input{\subfix{metadata.tex}}
\PaperHeading

\section{Introduction \& Motivation}
Estimating depth from a single camera image is useful wherever a depth sensor
is too expensive or too large, for example in small robots, driver assistance
or augmented reality. Training such a network requires a depth value for
every pixel, which is costly to measure in real scenes. Game engines and
renderers, in contrast, store the exact distance of every pixel in their depth
buffer, so synthetic images come with perfect labels for free. The catch is
that a network may learn more than depth: it can also learn regularities that
only hold in its virtual training world, so-called shortcuts
\autocite{gamedepth-geirhos2020}, and these may fail on real images.

This project asks: \emph{What does a small convolutional network learn about
depth from game data alone, and which shortcuts of the virtual world does it
carry over to real webcam images?} The question is made concrete by the live
demo at the KI-CON: a visitor holds a hand in front of the body, and the depth
map should show the hand closer than the torso. Telling a hand in front of the
body from a hand behind it is a small, checkable task that average error
metrics hardly capture. Beyond the demo, the answer matters to everyone who
trains models on synthetic data for cheap hardware, for instance for
low-cost robots or assistive devices, because a shortcut that goes unnoticed
in the virtual test set can fail silently in use.

The contribution is threefold: (i) a compact U-Net trained only on rendered
data that runs in real time on a laptop; (ii) an evaluation beyond average
errors, with a hand-order test, a flat baseline, a height test and a fair
comparison with Depth Anything~V2; and (iii) a webcam test with an unambiguous
protocol and paired statistics that shows which shortcuts transfer to real
images. The goal is not to beat large foundation models but to understand a
small, cheap one.

\section{State of the Art}
Eigen et al.\ trained the first convolutional network for single-image depth
and introduced a scale-invariant loss, because absolute scale is ambiguous in
a single image \autocite{gamedepth-eigen2014}. Encoder--decoder networks with
skip connections such as the U-Net \autocite{gamedepth-ronneberger2015} and
ImageNet-pretrained encoders such as ResNet \autocite{gamedepth-he2016} have
since become standard building blocks.

Current models generalise by scale. MiDaS mixes many datasets and evaluates
with an affine alignment of inverse depth \autocite{gamedepth-ranftl2022}.
Depth Anything~V2 trains a teacher on 595{,}000 synthetic images and students
on more than 62~million pseudo-labelled real images; its smallest version
still has 24.8\,M parameters \autocite{gamedepth-yang2024}. Depth Pro predicts
sharp metric depth \autocite{gamedepth-bochkovskii2025}, and Depth Pro~2,
published in November~2026, extends this to temporally consistent video with
models from 24\,M to 872\,M parameters \autocite{gamedepth-hu2026}. These
models transfer well, but they need data and compute far beyond a student
project.

Games have been used as a source of ground truth before
\autocite{gamedepth-richter2016}; MVS-Synth provides depth from GTA~V
\autocite{gamedepth-huang2018} and RHD renders people with hands close to the
camera in Blender \autocite{gamedepth-zimmermann2017}. Van Dijk and de Croon
showed that depth networks rely on the vertical image position of objects,
since in driving scenes ``lower in the image'' usually means ``closer''
\autocite{gamedepth-vandijk2019}. Little is known about what a small network
trained \emph{only} on such data learns for a close-range task like the hand
in front of the body, and whether its shortcuts survive on real images. This
is the gap addressed here.

\section{Method}
\begin{figure}[t]
  \centering
  \includegraphics[width=\columnwidth]{\subfix{figures/gamedepth-pipeline.pdf}}
  \caption{Overview: training only on rendered data, evaluation on rendered
  and real images, and the live demo.}
  \label{fig:gamedepth-pipeline}
\end{figure}
Figure~\ref{fig:gamedepth-pipeline} summarises the approach.
\textbf{Data.} The main dataset is RHD \autocite{gamedepth-zimmermann2017}:
41{,}258 training and 2{,}728 test images ($320\times320$) of rendered people
with depth and masks for body and hand. It was chosen because, unlike GTA,
it shows hands close to the camera, which is the demo situation. A data check
showed two properties that matter later: the background is a flat photograph
at a constant 5\,m and covers more than half of all pixels, and the hand is in
front of the body in 88.6\,\% of the test images (base rate). As a contrast we
use 11{,}000 GTA~V street scenes from MVS-Synth \autocite{gamedepth-huang2018};
their depth is given in game units, and values of 1000 and more are treated
as invalid (9.5\,\% of pixels).

\textbf{Network.} A U-Net \autocite{gamedepth-ronneberger2015} with an
ImageNet-pretrained ResNet-18 encoder \autocite{gamedepth-he2016}; each
decoder stage upsamples bilinearly, concatenates the skip features and applies
two $3\times3$ convolutions with batch normalisation and ReLU. The output is
log depth; the network has 14.4\,M parameters. We chose a convolutional
encoder over a vision transformer as used by Depth Anything because it is
small, fast at low resolution and works well with pretrained weights.

\textbf{Loss and training.} With $d_i=\log\hat y_i-\log y_i$ over the $n$
valid pixels, we minimise the scale-invariant loss \autocite{gamedepth-eigen2014}
\begin{equation}
  L_{\mathrm{si}}=\sqrt{\tfrac{1}{n}\textstyle\sum_i d_i^2-\lambda\big(\tfrac{1}{n}\sum_i d_i\big)^2},\quad \lambda=0.85,
  \label{eq:gamedepth-loss}
\end{equation}
plus 0.5 times an L1 term on the horizontal and vertical differences of $d$ at
three scales, which keeps edges sharp. Scale invariance suits the demo, where
only the order of distances matters, and lets us mix metres (RHD) with game
units (GTA). Training uses random $256\times256$ crops, horizontal flips and
colour augmentation (brightness, contrast, saturation, colour cast, blur,
noise), AdamW with a one-cycle schedule (peak learning rate $3\cdot10^{-4}$),
batch size 32, 15 epochs and mixed precision on a Colab L4 GPU (61\,min for
RHD). We trained an RHD network, a GTA network and a mixture of 70\,\% RHD and
30\,\% GTA drawn by a weighted sampler. We report the last epoch, since picking
the best epoch would select on the test set.

\textbf{Evaluation on rendered data.} On the 2{,}728 RHD test images
(centre crops) we measure AbsRel and $\delta_1$ \emph{on the person only},
after the same affine alignment in inverse depth for every model (MiDaS
protocol, clamped at 5\,m), because the Depth Anything~V2 model used here
predicts relative inverse depth. A flat prediction with one depth everywhere serves as baseline.
The \emph{hand test} compares the median predicted depth in the hand mask with
that in the body mask. Because a model that always answers ``hand in front''
already reaches the base rate of 88.6\,\%, we report the balanced accuracy,
the mean of the hit rates for ``in front'' and ``behind''; guessing gives 0.5.
The \emph{height test} splits images by whether the hand is lower in the image
than the body and measures how much the share of ``in front'' answers changes
with the true order (influence of truth) versus with the image height
(influence of height).

\textbf{Webcam test.} A recording tool takes webcam photos and lets the user
draw boxes around hand and torso and label the order; it deliberately hides
the depth map so that poses are not chosen in favour of a model. A model's
answer compares the median predicted nearness in the central 60\,\% of both
boxes. A pilot with 32 photos showed that ``hand behind the head'' was
ambiguous. In the final session (41 photos) the person stood 40--50\,cm in
front of a wall: ``behind'' means the hand lies flat on the wall with the arm
stretched, next to the head or the hip; ``in front'' means 30--40\,cm in
front of the chest at head or belly height. We give 95\,\% bootstrap intervals (2{,}000 resamples) and
compare models on the same photos with an exact McNemar test
\autocite{gamedepth-mcnemar1947}.

\section{Results \& Evaluation}
\textbf{Rendered test set.} Table~\ref{tab:gamedepth-rhd} shows that the RHD
network estimates the depth of persons best, 35\,\% below the flat baseline
and below Depth Anything~V2~Small. In ordering the hand it is weaker than
Depth Anything (balanced accuracy 0.83 vs.\ 0.90): it detects 95\,\% of hands
in front but only 71\,\% of hands behind, whereas Depth Anything reaches 81\,\%
and 98\,\%. The GTA network is at chance level, and its answers depend more on
image height (influence 0.28) than on the true order (0.07), as predicted by
van Dijk and de Croon \autocite{gamedepth-vandijk2019}. Figure~\ref{fig:gamedepth-rhd}
illustrates this: it renders the raised hand darker, that is, farther away.
The mixed network does not take over this shortcut on rendered data, and its
depth error stays within the run-to-run variation; two runs with identical settings
differed by 0.025 in whole-image AbsRel. ImageNet pretraining is the clearest
design effect (whole-image AbsRel 0.13--0.16 with vs.\ 0.23--0.24 without);
augmentation only shows a tendency.

\begin{table}[t]
  \centering
  \caption{Rendered RHD test set (2{,}728 images). AbsRel and $\delta_1$ on
  the person after identical affine alignment; BA = balanced accuracy of the
  hand test; height = influence of image height.}
  \label{tab:gamedepth-rhd}
  \small\setlength{\tabcolsep}{4pt}
  \begin{tabular}{lcccc}\toprule
    Model & AbsRel$\downarrow$ & $\delta_1\uparrow$ & BA$\uparrow$ & height \\ \midrule
    Flat baseline      & 0.176 & 0.751 & 0.50 & 0.00 \\
    GTA net            & 0.164 & 0.783 & 0.51 & 0.28 \\
    Mix 70/30          & 0.119 & 0.891 & 0.80 & $-0.10$ \\
    RHD net            & \textbf{0.114} & \textbf{0.903} & 0.83 & $-0.07$ \\
    Depth Anything V2-S & 0.132 & 0.877 & \textbf{0.90} & $-0.01$ \\ \bottomrule
  \end{tabular}
\end{table}

\begin{figure}[t]
  \centering
  \includegraphics[width=\columnwidth]{\subfix{figures/gamedepth-rhd-examples.pdf}}
  \caption{Rendered test images; bright means near. The RHD net follows the
  ground truth, the GTA net renders the raised hand dark (far).}
  \label{fig:gamedepth-rhd}
\end{figure}

\textbf{Real webcam photos.} Table~\ref{tab:gamedepth-webcam} reports the
final session. Depth Anything is correct on all 41 photos, which confirms that
the labels are unambiguous; in the pilot it had recognised only 5 of 15 hands
behind the body. Our networks recognise every hand in front of the body. The
RHD network also finds 17 of 25 hands behind it, so it does not simply answer
``in front'', which would yield 0 of 25. Its errors concentrate on one pose:
the hand behind the body at hip height. The paired comparison confirms that
Depth Anything is better on this test (8\,:\,0 discordant photos against the
RHD net, $p=0.008$; 11\,:\,0 against the mix, $p=0.001$), whereas the
difference between the RHD and the mixed network is not significant
(5\,:\,2, $p=0.45$). On real photos the height influence rises to 0.31 for the
RHD and 0.42 for the mixed network, although it was close to zero on rendered
data.

\begin{table}[t]
  \centering
  \caption{Webcam photos, final session: correctly ordered photos per pose
  (high/low: hand box centre above/below the torso box) and balanced accuracy
  (BA) with 95\,\% bootstrap interval.}
  \label{tab:gamedepth-webcam}
  \small
  \begin{tabular}{lccc}\toprule
    Hand & RHD & Mix & DA~V2-S \\ \midrule
    in front, high & 8/8 & 8/8 & 8/8 \\
    in front, low  & 8/8 & 8/8 & 8/8 \\
    behind, high   & 12/12 & 12/12 & 12/12 \\
    behind, low    & 5/13 & 2/13 & 13/13 \\ \midrule
    BA             & 0.84 & 0.78 & 1.00 \\
    95\,\% interval & {\scriptsize 0.74--0.93} & {\scriptsize 0.68--0.88} & {\scriptsize 1.00--1.00} \\ \bottomrule
  \end{tabular}
\end{table}

\begin{figure}[t]
  \centering
  \includegraphics[width=\columnwidth]{\subfix{figures/gamedepth-webcam-failure.pdf}}
  \caption{Typical failure (cropped below the face): the hand lies flat on
  the wall behind the person. Depth Anything shows it as far away as the wall,
  both own networks as near as the hip.}
  \label{fig:gamedepth-fail}
\end{figure}

\textbf{Speed.} On a laptop with an RTX~4050 the RHD network needs 5.8\,ms per
image at $320\times256$ (32-bit), Depth Anything~V2~Small 27.5\,ms at
$518\times518$ (22.1\,ms in 16-bit); on the CPU the times are 103 and 628\,ms.
The built-in webcam itself delivers at most about 20 frames per second,
which limits the live demo.

\section{Discussion \& Limitations}
The results support the research question in two directions. The network
learned real depth cues: it estimates the depth of persons well and transfers
the main demo case, a hand in front of the body, to real photos without error. It
also learned three shortcuts of its rendered world. First, the base rate makes
it favour ``hand in front''. Second, because the RHD background is always
5\,m away, it treats everything except the person as far; real rooms appear
uniformly dark. Third, Figure~\ref{fig:gamedepth-fail} suggests that it
assigns the hand the depth of the body part at the same height: next to the
hip, which is closer to the low webcam than the chest, the hand looks ``in
front''. This third explanation is a hypothesis; a counter-test would train
with images flipped upside down at random, which removes the height cue. The shortcuts
appeared only partly on rendered data, so a real test set was essential.

Compared with Depth Anything~V2~Small, the network trades reliability for
efficiency: it uses about 1{,}500 times fewer training images, 1.7 times fewer
parameters, and is 4.7 times faster on the GPU (3.8 times against 16-bit
Depth Anything) and 6 times on the CPU. Part of the speed gain comes from the
smaller input. Whether it runs on smaller devices was not tested.

The evidence is limited: one person, one room and one camera, so effects of
skin tone, clothing or lighting are untested; boxes instead of pixel masks; 41 photos with wide intervals; one or two training runs per
variant; and only three mixing ratios (0, 30, 100\,\% GTA). Training data with
structured backgrounds and hands far behind the body, more participants and a
comparison with newer models such as Depth Pro~2 are concrete next steps.
The webcam photos show only the author and are not part of the public
repository.

\section{Conclusion}
A 14.4\,M-parameter U-Net trained only on rendered data estimates the depth of
persons well, runs in real time and handles the hand in front of the body on
real webcam photos. A targeted test shows, however, that it also learned the
base rate, the flat background and probably a height cue of its virtual world,
which Depth Anything~V2 avoids in our tests. A next step is the flip
counter-test to confirm the height explanation. \emph{Use of AI tools:} code
and text drafts were written with the help of an AI assistant (Claude,
Anthropic); the author ran and checked all experiments and results.

\printbibliography[heading=subbibliography,title={References}]
\end{refsection}
\end{document}
